\section{Layered histories and the analytic inputs}\label{sec:dhm-inputs}

We use the version of Deng--Hani--Ma dated 18 July 2025,
\cite{DHM2025}. All source page and equation numbers in this section
refer to that version. Its diagram estimates concern positive geometric
integrals. They do not assert a fluctuation theorem or the sharp
connected estimate proved below.

\subsection{Parameters and the restarted dynamics}

Fix \(\upsilon=3^{-4}=1/81\). The constant \(\Gamma\), which bounds
collision excess in one collision cluster, is chosen as a sufficiently
large dimensional constant before the number of slabs \(L\).
Put \(\tau=T/L\), and, with harmless integer roundings, set
\begin{equation}\label{eq:dhm-cutoffs}
 A_L=|\log\eps|,\qquad
 \Lambda_i=A_i^{30},\qquad A_{i-1}=\Lambda_i^{30}.
\end{equation}
For the analytic induction use
\begin{equation}\label{eq:dhm-exponents}
 \beta_i=\Bigl(1-\frac{i}{10L}\Bigr)\frac\beta2,
 \qquad \widetilde\beta_i=\frac{\beta_{i-1}+\beta_i}{2},
 \qquad
 \theta_i=\Bigl(1-\frac{i}{10L}\Bigr)3^{-11}.
\end{equation}
Constants denoted by \(C\) in an exponential factor \(C^b\) can be
chosen independently of \(L\); constants \(C_*\), and positive constants
\(c_*\), can depend on \(L\) and the kinetic bounds. All are independent
of \(\eps\). No moment-order threshold is used in choosing
\(\Gamma,L\).

The restarted truncated dynamics is defined on all configurations, with
overlaps permitted. Within slab \(i\), start the collision prehistory
anew. At each incoming contact, accept the elastic collision precisely
when adding that contact to the existing collision graph leaves its
component with at most \(\Lambda_i\) particle labels and excess at most
\(\Gamma\). Otherwise both particles continue straight. A contact
inside one existing component increases its excess by one; a contact
between different components merges them and adds their excesses.
Here, for a graph with \(b\) edges, \(n\) vertices, and \(k\) components,
its excess is \(r=b-n+k\), counting repeated edges with multiplicity.
The caps bound the number of accepted collisions in a finite system;
between accepted collisions the motion consists of finitely many
straight segments. Incoming contacts, rather than exits from an
overlap, determine the rule. Binary changes of variables exclude the
null sets of grazing contacts, simultaneous contacts, and branch ties.
This is the rule of \cite[Definition 4.3, Equation (4.2)]{DHM2025}.
Almost-everywhere existence, the cluster caps, and mass preservation are
\cite[Proposition 4.5]{DHM2025}. Restarting is part of our definition;
the capped evolution with a prehistory retained is not a semigroup.
We call the evolution obtained by concatenating these restarted slab
dynamics the pasted dynamics.

\begin{lemma}[Kinetic envelopes]\label{lem:kinetic-envelopes}
Under \Cref{ass:kinetic},
\begin{equation}\label{eq:kinetic-bol-propagated}
 \sup_{0\le t\le T}
 \bigl(\|f_t\|_{\Bol,\beta}
       +\|\nabla_x f_t\|_{\Bol,\beta}\bigr)<\infty.
\end{equation}
\end{lemma}
\begin{proof}
Apply \cite[Proposition A.1]{DHM2025} with decreasing Gaussian exponent
\(\beta_*(t)=2\beta-\beta t/T\). The weighted uniform hypothesis follows
from \eqref{eq:kinetic-gaussian}, and its two initial summable spatial
norms are \eqref{eq:initial-bol-bounds}. Since
\(\beta_*(t)\ge\beta\), its conclusion implies
\eqref{eq:kinetic-bol-propagated}. This is the same application as
\cite[Equation (1.29)]{DHM2025}.
\end{proof}

\subsection{Histories, roots, and normalization}

\begin{definition}[Layered C/O history]\label{def:dhm-history}
A history consists of finitely many oriented particle lines and binary
atoms, with an order on the atoms along every line. Each atom has two
incoming and two outgoing half-edges, paired into two serial particle
lines. A C-atom applies elastic scattering, and an O-atom records an
incoming encounter with straight continuation. A line segment between
atoms is a bond. Segments with only one incident atom are ends. A line
with no atom is a single empty end; it is integrated unless its state
is fixed, just like any other end. Atoms belong to specified slabs; the order
along every line must be consistent with those slabs. Lines can begin
or end at slab interfaces. Their beginnings carry one-particle input
functions. Certain top lines, with distinct prescribed labels, are
roots. The remaining particle names are immaterial: histories are
identified under relabelings preserving the roots, the serial pairing,
the C/O types, and the line orders. Integrals use the factorial weights
obtained from ordered particle labels, as specified in
\Cref{sec:cluster-expansion}.

In a slab, the C-clusters are the components formed using only C-atoms.
The cluster graph has those clusters as vertices and the O-atoms as
edges. Its components are called layer components. Isolated rooted
lines are included. Let \(H_i\) be the roots of slab \(i\): the lines
continued into the part of the history above that slab, or the designated
roots if it is the highest slab. Let \(H_0\) be the lines incident to at
least one initial link. Initial links join pairs of lines born at zero
and carry the positive weight
\begin{equation}\label{eq:dhm-initial-link}
 \ell_\eps(x,x')
 =\ind_{\{|x-x'|\le\eps\}}
   +\eps^\upsilon\ind_{\{|x-x'|>\eps\}}.
\end{equation}
An admissible dependent history satisfies:
\begin{enumerate}[label=(\roman*)]
 \item \(|H_i|\le A_i\); every C-cluster has at most \(\Lambda_i\)
 particle lines and excess at most \(\Gamma\); every layer component
 contains fewer than \(\Lambda_i\) C-clusters.
 \item The cluster graph in each slab is a forest, and every layer
 component contains a root of that slab. Every unrooted line is
 connected in that slab to a rooted line.
 \item The initial links form a forest.
 \item For \(i\ge2\), each line in \(H_i\) is connected within slab
 \(i\) to a different line in \(H_i\), connected there to a line in
 \(H_{i-1}\), or itself belongs to \(H_{i-1}\). For \(i=1\), replace
 \(H_{i-1}\) by \(H_0\).
\end{enumerate}
Connection within a slab means intersection with the same component of
its atom graph; the last alternative in (iv) covers a straight line
with no atom there. These are the conditions of
\cite[Definition 3.20(1)]{DHM2025}. Writing \(\mathcal M_i\) for
the history restricted to slab \(i\), its complexity is
\begin{equation}\label{eq:dhm-complexity}
 \rho=|H_0|+\sum_i\bigl(|H_i|+r(\mathcal M_i)\bigr),
\end{equation}
where \(r(\mathcal M_i)\) is its total cycle excess. Since the
O-graph is a forest, this equals the sum of the C-cluster excesses.
The total number of atoms is denoted by \(b\).
\end{definition}

Here is an explicit integral convention, also needed after cutting a
history. Give each edge \(e\) a vector \(z_e=(x_e,v_e)\), where its
physical position at time \(t\) is \(X_e(t)=x_e+tv_e\). Give atom \(a\)
a time \(t_a\). Its two incoming edges are \(1,2\), its serial outgoing
edges \(1',2'\). On its contact surface put
\(u=v_1-v_2\), \(\omega=(X_1(t_a)-X_2(t_a))/\eps\), and
\(\mathfrak b=[-u\cdot\omega]_+\). Define
\begin{align}
 \Delta_a^{\mathrm C}
 &=\delta(|X_1(t_a)-X_2(t_a)|-\eps)\mathfrak b
   \prod_{j=1}^2\delta(X_{j'}(t_a)-X_j(t_a))
   \delta(v_{1'}-v_1+(u\cdot\omega)\omega)
   \delta(v_{2'}-v_2-(u\cdot\omega)\omega),\label{eq:atom-C}\\
 \Delta_a^{\mathrm O}
 &=\delta(|X_1(t_a)-X_2(t_a)|-\eps)\mathfrak b
   \prod_{j=1}^2\delta(x_{j'}-x_j)\delta(v_{j'}-v_j).
   \label{eq:atom-O}
\end{align}
Vector deltas in these formulas are three-dimensional. Reversing
\(\omega\) converts the incoming convention to
\eqref{eq:collision-kernel}; surface measure is unnormalized.
For a cut history some ends are fixed. If \(E_*\) consists of all bonds
and nonfixed ends, set
\begin{equation}\label{eq:edge-integral}
 I_M(Q)(z_{\mathrm{fixed}})
 :=\mu^{|E_*|-2b}
   \int_{\mathcal D_M}\int
    Q(z_E,t_M)\prod_{a\in M}\Delta_a
       \prod_{e\in E_*}\dd z_e\prod_{a\in M}\dd t_a.
\end{equation}
The domain \(\mathcal D_M\) imposes the assigned slab on each time and
strictly increasing times along every particle line. These positive
measures can equivalently be defined by the incoming surface change of
variables, so the displayed deltas are a notation for the corresponding
coarea measures. The kernel \(Q\) contains the birth factors evaluated
at physical positions, the initial-link factors, and any further
history indicators. All statements below allow such indicators.

For a full history, \(|E|-2b=n\), the number of particle lines. If there
are \(h\) designated roots, its normalized positive integral is
\begin{equation}\label{eq:normalized-history}
 \mathcal J_M(Q)=\mu^{-h}I_M(Q).
\end{equation}
Thus every birth has activity \(\mu\), every root has normalization
\(\mu^{-1}\), and every contact surface supplies a factor \(\eps^2\).
This is \cite[Proposition 7.5, Equation (7.7)]{DHM2025}. To keep a root
endpoint \((x,v)\) at time \(t\) unintegrated, turn its top end into a
fixed end and evaluate at \((x-tv,v)\). This removes the root activity
factor from \eqref{eq:edge-integral}; its resulting density has the
normalization of \eqref{eq:normalized-history}.

\subsection{The imported analytic package}

The one-slab class \(\mathcal T_i\) consists of connected histories
with one top root, the size, excess, and forest conditions (i)--(ii),
and no continuation or initial-link requirement. In particular,
\(\mathcal T_i\) can have order \(\Lambda_i^2\) lines. Let
\(\mathcal T_i^{\mathrm{err}}\) be the histories obtained by inserting
one additional O-atom in such a history, with a compatible position
in each of the two line orders. The leading signed sum contains every
one-root particle-line tree of total size \(<\Lambda_i/3\), with its
C/O choices and sign \((-1)^{\#\mathrm O}\), once with the factorial
weight coming from particle labels. It does not impose the absence of
additional encounters.

\begin{theorem}[DHM analytic package]\label{thm:dhm-package}
Assume \Cref{ass:kinetic}, the parameter hierarchy
\eqref{eq:dhm-cutoffs}--\eqref{eq:dhm-exponents}, and a common bound on
the Gaussian Bol norms of all birth inputs. For sufficiently large
fixed \(L\) and sufficiently small \(\eps\), the following hold.
\begin{enumerate}[label=(\alph*)]
 \item For admissible dependent histories with prescribed labeled roots,
 fixed \(b\), and complexity \(\rho\), the sum of their normalized
 positive integrals is at most
 \begin{equation}\label{eq:dhm-coarse}
 C^b\tau^{b/10}\eps^{c_*\rho}
       |\log\eps|^{C_*\rho}.
 \end{equation}
 The sum includes the graph enumeration and the admissible indicator
 restrictions; a bounded observable multiplier costs its supremum.
 The birth inputs need not equal the approximate density constructed
 in the source, but must satisfy its common Gaussian Bol envelope.

 \item Suppose the independent input \(g\) in slab \(i\) satisfies
 \begin{equation}\label{eq:dhm-input-closeness}
 \|g-f_{(i-1)\tau}\|_{\Bol,\beta_{i-1}}
       \le\eps^{\theta_{i-1}}.
 \end{equation}
 Let \(\mathsf T_i^\eps[g]\) be the leading signed sum just defined.
 Then
 \begin{equation}\label{eq:dhm-signed-trees}
 \|\mathsf T_i^\eps[g]-f_{i\tau}\|_{\Bol,\beta_i}
       \le\tfrac12\eps^{\theta_i}.
 \end{equation}
 In the fixed-top-endpoint norm \(\Bol,\widetilde\beta_i\), the summed
 absolute contributions of cyclic members of \(\mathcal T_i\), members
 of \(\mathcal T_i^{\mathrm{err}}\), and tree members of
 \(\mathcal T_i\) with at least \(\Lambda_i/3\) lines are
 \(o(\eps^{\theta_i})\). Fixed powers of the cutoffs, arising from a
 bounded number of selected edge positions, can be included. A fixed
 exponential factor per atom can also be included by choosing \(L\)
 sufficiently large.
\end{enumerate}
Part (b) also holds uniformly when the final slab stops at any
\(0\le\sigma\le\tau\). With the same cutoffs and exponent budgets,
write \(\mathsf T_{i,\sigma}^\eps\) for its signed tree sum. Then
\begin{equation}\label{eq:dhm-partial-slab}
 \sup_{0\le\sigma\le\tau}
 \|\mathsf T_{i,\sigma}^\eps[g]
          -f_{(i-1)\tau+\sigma}\|_{\Bol,\beta_i}
       \le\tfrac12\eps^{\theta_i},
\end{equation}
and the positive error classes obey the same full-\(\tau\) upper
majorants uniformly in \(\sigma\). At \(\sigma=0\) the map is the
identity and its contact error classes vanish.
The assertion in (b) is about the specified complete signed tree sum
and positive error classes. Identifying a different one-root operator
with that sum up to those errors requires a separate combinatorial
argument.
\end{theorem}

\begin{proof}[Source statements and the envelope adaptations]
For (a), \cite[Proposition 7.2]{DHM2025} counts the root-labeled
unlabeled histories by \(C^b|\log\eps|^{C_*\rho}\).
\cite[Proposition 9.7, Equation (9.51)]{DHM2025} retains a positive
geometric gain proportional to \(\rho\). After the regular cutting
operations, \cite[Equation (9.53)]{DHM2025} bounds each full edge
integral by
\[
 \tau^{b/9}\eps^{-2h+2\upsilon/5+c_*\rho}
       |\log\eps|^{C_*\rho}.
\]
Multiplication by \(\mu^{-h}=\eps^{2h}\) cancels the negative root
power. Incorporating the counting and the cutting alternatives,
weakening \(1/9\) to \(1/10\), and discarding the extra positive power
gives \eqref{eq:dhm-coarse}. These are the estimates in the proof of
\cite[Proposition 6.2, Section 9.2]{DHM2025}, rather than just its final
displayed conclusion.

The only use there of a birth density is the decomposition in
\cite[Equations (9.58)--(9.63)]{DHM2025}: write its absolute value as a
sum of spatially localized Gaussian envelopes, whose coefficients have
bounded sum. Indeed, from \eqref{eq:bol-norm} choose a summable sequence
\((m_a)\) bounding the weighted supremum on the unit balls; the bounded
overlap of those balls gives a decomposition dominated by
\(C\sum_a m_a\ind_{|x-a|\le1}e^{-\beta'|v|^2}\), at any permitted
exponent \(\beta'>0\). Apply the positive source estimates to each
product of such envelopes and sum their coefficients. This proves the
stated extension to arbitrary bounded inputs, including absolute values
of signed inputs. Multipliers between zero and one are kept in the
support throughout, or dropped only when a positive upper bound is
being taken.

For (b), the one-root integral is a function of a fixed endpoint by
\cite[Equation (14.2)]{DHM2025}. The C/O involution of
\cite[Proposition 14.2]{DHM2025} cancels the noncausal trees in the
complete signed sum. The remaining causal trees satisfy the Gaussian
Bol estimate \cite[Proposition 14.3, Equation (14.8)]{DHM2025}.
The comparison to the short-time Boltzmann expansion, using
\Cref{lem:kinetic-envelopes}, is the proof of
\cite[Proposition 6.1, Section 14]{DHM2025}. To replace its preceding
input by \(g\), expand each birth factor as
\(f_{(i-1)\tau}+(g-f_{(i-1)\tau})\). Each residual factor gains
\(\eps^{\theta_{i-1}}\) by \eqref{eq:dhm-input-closeness}; the source
multilinear tree bound sums all choices. Since
\(\theta_{i-1}>\theta_i\), the residual and the finite-diameter
comparison errors are \(o(\eps^{\theta_i})\), proving
\eqref{eq:dhm-signed-trees} for small \(\eps\).

For the positive subleading terms the explicit source bound is
\cite[Equation (14.28)]{DHM2025}, in
\(\Bol,\widetilde\beta_i\), with a time power per atom and gain
\(\eps^{\upsilon/3+c_*\rho}\) in the cyclic and extra-encounter cases.
Points 1--2 of Section 14.3 retain the Gaussian weight at the fixed
root and the spatial Bol sum; Point 4 regularizes an O-atom adjacent
to that fixed root without altering its integral. Thus this is stronger
than an integrated mass estimate. The tree-tail branch of Point 5 uses
only that the atom number is of cutoff order; replacing
\(\Lambda_i\) by \(\Lambda_i/3\) leaves an exponentially small tail.
Graph enumeration, the bounded number of selected-edge choices, and
fixed exponential factors are absorbed by the time power after taking
\(L\) large. This proves (b) with the asserted norm and size range.
For the partial-slab assertion, restrict every contact time to
\(((i-1)\tau,(i-1)\tau+\sigma)\). After translating the slab to start
at zero, the fixed top end is evaluated at \((x-\sigma v,v)\).
Repeat the root Gaussian and spatial localization argument of
\cite[Section 14.3, Points 1--2]{DHM2025} at this endpoint: all spatial
shifts have duration at most \(\tau\), and the same Gaussian exponent
margin remains available. The resulting fixed-root Bol bounds are
therefore uniform for \(\sigma\le\tau\). Positivity enlarges the
contact-time domain only with the intercept held fixed; the preceding
localization supplies the bound in the physical endpoint norm.
The signed tree involution
preserves this restriction, and the Boltzmann expansion and its
finite-diameter comparison use the same shortened time domain.
Their upper estimates retain the original \(\tau\) bound, uniformly
in \(\sigma\), so the preceding argument proves
\eqref{eq:dhm-partial-slab} with unchanged budgets. When
\(\sigma=0\), only the zero-atom tree remains, giving \(g\);
\eqref{eq:dhm-input-closeness} and
\(\theta_{i-1}>\theta_i\) give the claimed bound there too.
\end{proof}

\begin{lemma}[Coarse boundary modifications]\label{lem:dhm-dummies}
The following modifications of the positive coarse bound are allowed.
\begin{enumerate}[label=(\roman*)]
 \item If continuation condition (iv) is omitted in the top slab with
 \(h\) roots, the bound may lose \(\eps^{-C_*h}\).
 \item A connected pure-C terminal component with at most
 \(3\Lambda_i\) lines and excess at most \(2\Gamma\), meeting a
 continuation below but having no designated top root, can be treated
 by adding one artificial root at cost \(\mu\).
 \item A fixed number \(q\) of specified lines, continued to zero, can
 be supplied with initial-link continuations by adding \(q\) dummy
 straight root lines at cost at most \(\eps^{-q\upsilon}\), provided
 the root cutoffs have the corresponding fixed slack.
\end{enumerate}
Bounded records carried by these lines can be kept in the integral.
\end{lemma}
\begin{proof}
For (i), use the error version of the cutting inequality in
\cite[Equation (9.51)]{DHM2025}, retaining its deficit proportional to
\(h\). This is an adaptation of the source proof: its formal error
class also requires a large terminal component. The finale of
\cite[Section 13.5, p.~169]{DHM2025} treats the missing top continuation
in two cases. When \(\rho\) is at most a fixed multiple of \(C_*h\),
the source's UP cutting algorithm produces at most \(h\) elementary
pieces carrying a free-component activity cost, called bad pieces in
that proof. The permitted root deficit pays for them. For larger
\(\rho\), the same cutting argument as for dependent histories applies;
the continuation condition used in its layer selection concerns a
strictly lower layer, where that condition is still available. Thus
the cutting inequality with its deficit does not require a large
terminal component. The subsequent positive integrations give the loss
\(\eps^{-C_*h}\) in \cite[Equation (9.82)]{DHM2025}. Terminal size is
used only later, when the time tail absorbs that deficit. Keeping the
deficit gives precisely (i).

For (ii), choose one line of the terminal component as root. Its
normalization multiplies the original integral by \(\mu^{-1}\);
multiplication by \(\mu\) restores it. The stated enlarged terminal
class must be split into two cases for the source application.
If the component has at most \(\Lambda_i\) lines and excess at most
\(\Gamma\), its root and the stated lower continuation make it an
ordinary admissible dependent history.  If either cutoff is exceeded,
it belongs to the enlarged terminal class in
\cite[Definition 3.20(4)]{DHM2025}, and the same root operation is
\cite[Equation (15.4)]{DHM2025}. A sum over the choice of line costs
a cutoff polynomial.

For (iii), take any fixed nonnegative \(q_0\) with
\(\int q_0\,\dd z=1\) and finite Gaussian Bol norm. For every \(x\),
\[
 \int q_0(x',v')\ell_\eps(x,x')\,\dd x'\dd v'
       \ge\eps^\upsilon.
\]
Insert this integral once for each required continuation. The activity
\(\mu\) of a dummy input cancels its additional root factor
\(\mu^{-1}\). Its straight line satisfies continuation in every slab,
and the new links are pendant edges, preserving the initial forest.
The positive kernel inequality proves the claimed cost independently
of whether \(f_0\) vanishes anywhere. The envelope extension in
\Cref{thm:dhm-package} permits \(q_0\) as the dummy input. Added roots
and the finitely many existing roots used for their records fit the
prescribed slack. No record is changed by this multiplication.
\end{proof}
